% 贯穿图的子图与诱导子图：S={1,2,3}。虚线灰色表示被删除的顶点或边。
\begin{tikzpicture}[panelmark/.style={font=\figurefont\fontsize{8}{10}\selectfont,text=FigureInk,anchor=north,inner sep=0pt,outer sep=0pt},
>=Stealth,
  v/.style={circle,draw=FigureStroke,fill=FigureBlueFill,line width=.8pt,minimum size=6.5mm,inner sep=0pt},
  vd/.style={circle,draw=FigureMuted,dashed,fill=white,text=FigureMuted,line width=.6pt,minimum size=6.5mm,inner sep=0pt},
  e/.style={draw=FigureStroke,line width=.8pt},
  ed/.style={draw=FigureMuted,dashed,line width=.6pt}]
\newcommand{\pos}{\coordinate (p1) at (-1.0,0); \coordinate (p2) at (0,.8);
  \coordinate (p3) at (0,-.8); \coordinate (p4) at (1.0,0);}
\begin{scope}[local bounding box=panelrow]
\begin{scope}[local bounding box=panela]
  \pos
  \node[v] (1) at (p1) {$1$}; \node[v] (2) at (p2) {$2$}; \node[v] (3) at (p3) {$3$}; \node[v] (4) at (p4) {$4$};
  \draw[e] (1)--(2) (1)--(3) (2)--(3) (2)--(4) (3)--(4);

\end{scope}
\begin{scope}[xshift=3.65cm,local bounding box=panelb]
  \pos
  \node[v] (1) at (p1) {$1$}; \node[v] (2) at (p2) {$2$}; \node[v] (3) at (p3) {$3$}; \node[vd] (4) at (p4) {$4$};
  \draw[e] (1)--(2) (1)--(3) (2)--(3); \draw[ed] (2)--(4) (3)--(4);

\end{scope}
\begin{scope}[xshift=7.3cm,local bounding box=panelc]
  \pos
  \node[v] (1) at (p1) {$1$}; \node[v] (2) at (p2) {$2$}; \node[v] (3) at (p3) {$3$}; \node[vd] (4) at (p4) {$4$};
  \draw[e] (1)--(2) (2)--(3); \draw[ed] (1)--(3) (2)--(4) (3)--(4);

\end{scope}
\begin{scope}[xshift=10.95cm,local bounding box=paneld]
  \pos
  \node[v] (1) at (p1) {$1$}; \node[v] (2) at (p2) {$2$}; \node[v] (3) at (p3) {$3$}; \node[v] (4) at (p4) {$4$};
  \draw[e] (1)--(2) (1)--(3) (2)--(4) (3)--(4); \draw[ed] (2)--(3);

\end{scope}
\end{scope}
% Subfigure marks: 8 pt, centred below each panel; common top 3 mm below this row.
\coordinate (panelrowbase) at ($(panelrow.south)+(0,-3mm)$);
\node[panelmark] at (panela.center |- panelrowbase) {(a)};
\node[panelmark] at (panelb.center |- panelrowbase) {(b)};
\node[panelmark] at (panelc.center |- panelrowbase) {(c)};
\node[panelmark] at (paneld.center |- panelrowbase) {(d)};
\end{tikzpicture}
